\documentclass[11pt]{article}

\usepackage{amsmath,amssymb,amsthm,mathtools}
\usepackage[margin=1in]{geometry}
\usepackage{microtype}

\newtheorem{theorem}{Theorem}[section]
\newtheorem{lemma}[theorem]{Lemma}
\newtheorem{proposition}[theorem]{Proposition}
\newtheorem{conjecture}[theorem]{Conjecture}
\newtheorem{remark}[theorem]{Remark}

\newcommand{\dd}{\,\mathrm d}
\newcommand{\one}{\mathbf 1}
\newcommand{\ip}[2]{\left\langle #1,#2\right\rangle}
\newcommand{\TT}{\mathsf T}

\title{Atlas 130: Exact Operator Reduction and Two Restricted-Family Theorems}
\author{}
\date{}

\begin{document}

\maketitle

\begin{abstract}
Atlas 130 is the graph consisting of two triangles joined by a bridge.
Its desired graphon inequality is
\[
 t(H_{130},W)\geq p^3(2p-1)^2\qquad (p\geq\tfrac12).
\]
We prove this inequality for two substantial classes: graphons whose
self-adjoint integral operator is positive semidefinite, and all regular
graphons.  For an arbitrary graphon we derive an exact one-line reduction
and isolate a stronger sufficient inequality as a conjecture, not a theorem.  An
exact rational two-step example shows why a tempting termwise proof of that
conjecture cannot work.  Thus this note gives rigorous restricted-family
progress but does not change the open status of Atlas 130.
\end{abstract}


%%%%%%%%%%%%%%%%%%%%%%%%%%%%%%%%%%%%%%%%%%%%%%%%%%%%%%%%%%%%%%%%%%%%%%%%%%%%%%%
\section{The graph and its target}
%%%%%%%%%%%%%%%%%%%%%%%%%%%%%%%%%%%%%%%%%%%%%%%%%%%%%%%%%%%%%%%%%%%%%%%%%%%%%%%

Use the Atlas labelling
\[
 E(H_{130})=\{01,02,03,12,34,35,45\}.
\]
Thus $012$ and $345$ are triangles and $03$ is their joining bridge.  The
Graph Atlas index is $130$ and the graph6 code is \verb|E{CW|.  Colouring
the first triangle, then the bridge endpoint, and then the last two triangle
vertices gives
\begin{equation}
 \chi_{H_{130}}(x)=x(x-1)^3(x-2)^2.
 \label{eq:chromatic}
\end{equation}
For $q=1-p$ and $r=2p-1$, the chromatic target is therefore
\begin{equation}
 q^6\chi_{H_{130}}(1/q)=p^3r^2.
 \label{eq:target}
\end{equation}
The formula on the left is interpreted polynomially at $p=1$.

Let $W\colon\Omega^2\to[0,1]$ be a symmetric measurable graphon on a
probability space.  Write
\[
 \TT g(x)=\int_\Omega W(x,y)g(y)\dd\mu(y),
 \qquad p=\int_{\Omega^2}W\dd\mu^2,
 \qquad d=\TT\one,
 \qquad A=\TT d.
\]
Define the rooted triangle density and three graph densities
\begin{align*}
 \tau(x)&=\int_{\Omega^2}
 W(x,y)W(x,z)W(y,z)\dd\mu(y)\dd\mu(z),\\
 B&=\ip dA=t(P_4,W),\\
 F&=\ip\tau A,\\
 H&=\ip\tau{\TT\tau}=t(H_{130},W).
\end{align*}
Here $F$ is the density of a triangle with a path of length two attached
at one triangle vertex.


%%%%%%%%%%%%%%%%%%%%%%%%%%%%%%%%%%%%%%%%%%%%%%%%%%%%%%%%%%%%%%%%%%%%%%%%%%%%%%%
\section{Two arbitrary-graphon anchor inequalities}
%%%%%%%%%%%%%%%%%%%%%%%%%%%%%%%%%%%%%%%%%%%%%%%%%%%%%%%%%%%%%%%%%%%%%%%%%%%%%%%

\begin{lemma}[Two-edge-tail anchor]
\label{lem:anchor}
For every graphon of density $p\geq\tfrac12$,
\begin{equation}
 F\geq rB
 \qquad\text{and}\qquad
 B\geq p^3.
 \label{eq:anchors}
\end{equation}
\end{lemma}

\begin{proof}
The elementary inequality $ab\geq a+b-1$, applied to
$W(x,y),W(x,z)$ and then multiplied by $W(y,z)$, gives
\begin{equation}
 \tau(x)\geq2A(x)-p.
 \label{eq:rooted-goodman}
\end{equation}
Since $A\geq0$ and $\int A=\lVert d\rVert_2^2$,
\[
 F\geq2\lVert A\rVert_2^2-p\lVert d\rVert_2^2.
\]
Put $g=d-p\one$ and $h=\TT g$, so $A=pd+h$.  A direct expansion and
completion of the square give
\begin{align}
 &2\lVert A\rVert_2^2-p\lVert d\rVert_2^2-r\ip Ad \notag\\
 &\quad=
 2\left\lVert h+\frac{2p+1}{4}g\right\rVert_2^2
 +\frac{(2p+1)(6p-1)}8\lVert g\rVert_2^2\geq0.
 \label{eq:anchor-square}
\end{align}
This proves $F\geq rB$ (indeed, the square works for $p\geq1/6$).

For completeness, the path bound follows from one application of
H\"older.  Define the quotients to be zero wherever their denominator
vanishes.  The integrand identity below holds almost everywhere because
zero-degree fibres carry no $W$-mass.  Thus
\begin{align*}
 p
 &=\int_{\Omega^2}
   \bigl(W(x,y)d(x)d(y)\bigr)^{1/3}
   \left(\frac{W(x,y)}{d(x)}\right)^{1/3}
   \left(\frac{W(x,y)}{d(y)}\right)^{1/3}\dd\mu^2\\
 &\leq B^{1/3}
 \left(\int\frac{W(x,y)}{d(x)}\dd\mu^2\right)^{1/3}
 \left(\int\frac{W(x,y)}{d(y)}\dd\mu^2\right)^{1/3}
 \leq B^{1/3}.
\end{align*}
Thus $B\geq p^3$.
\end{proof}


%%%%%%%%%%%%%%%%%%%%%%%%%%%%%%%%%%%%%%%%%%%%%%%%%%%%%%%%%%%%%%%%%%%%%%%%%%%%%%%
\section{The exact operator reduction}
%%%%%%%%%%%%%%%%%%%%%%%%%%%%%%%%%%%%%%%%%%%%%%%%%%%%%%%%%%%%%%%%%%%%%%%%%%%%%%%

\begin{proposition}[Only the negative spectral energy remains]
\label{prop:reduction}
Set
\[
 f=\tau-rd.
\]
Then the following are exact identities:
\begin{align}
 H-rF
 &=r(F-rB)+\ip f{\TT f},
 \label{eq:reduction-one}\\
 H-r^2B
 &=2r(F-rB)+\ip f{\TT f}.
 \label{eq:reduction-two}
\end{align}
\end{proposition}

\begin{proof}
Because $\tau=rd+f$ and $\TT d=A$,
\[
 H-rF=\ip\tau{\TT(\tau-rd)}
 =r\ip d{\TT f}+\ip f{\TT f}
 =r\ip A f+\ip f{\TT f}.
\]
But $\ip Af=F-rB$, proving \eqref{eq:reduction-one}.  Substituting
$F=rB+(F-rB)$ once more gives \eqref{eq:reduction-two}.
\end{proof}

Lemma~\ref{lem:anchor} makes $F-rB$ nonnegative.  Consequently the only
term in \eqref{eq:reduction-two} that can have the wrong sign for a general
graphon is the spectral energy $\ip f{\TT f}$.  A pointwise nonnegative
kernel need not define a positive-semidefinite operator, so this observation
does not by itself settle the arbitrary-graphon problem.


%%%%%%%%%%%%%%%%%%%%%%%%%%%%%%%%%%%%%%%%%%%%%%%%%%%%%%%%%%%%%%%%%%%%%%%%%%%%%%%
\section{Positive-semidefinite operator graphons}
%%%%%%%%%%%%%%%%%%%%%%%%%%%%%%%%%%%%%%%%%%%%%%%%%%%%%%%%%%%%%%%%%%%%%%%%%%%%%%%

\begin{theorem}[PSD-operator case]
\label{thm:psd}
Suppose the graphon operator is positive semidefinite on $L^2(\Omega)$:
\[
 \ip g{\TT g}\geq0\qquad\text{for every }g\in L^2(\Omega).
\]
Then, for every $p\geq\tfrac12$,
\[
 H\geq rF\geq r^2B\geq p^3r^2.
\]
In particular Atlas 130 satisfies its full chromatic target on this class.
\end{theorem}

\begin{proof}
In \eqref{eq:reduction-one}, all three quantities $r$, $F-rB$, and
$\ip f{\TT f}$ are nonnegative.  Thus $H\geq rF$.  The other two
inequalities are exactly Lemma~\ref{lem:anchor}.
\end{proof}


%%%%%%%%%%%%%%%%%%%%%%%%%%%%%%%%%%%%%%%%%%%%%%%%%%%%%%%%%%%%%%%%%%%%%%%%%%%%%%%
\section{All regular graphons}
%%%%%%%%%%%%%%%%%%%%%%%%%%%%%%%%%%%%%%%%%%%%%%%%%%%%%%%%%%%%%%%%%%%%%%%%%%%%%%%

\begin{theorem}[Regular case]
\label{thm:regular}
If $W$ is $p$-regular, meaning $d(x)=p$ almost everywhere, then for every
$p\geq\tfrac12$,
\[
 H\geq rF\geq p^3r^2.
\]
Thus Atlas 130 satisfies its full chromatic target for every regular
graphon, with no spectral assumption.
\end{theorem}

\begin{proof}
Let
\[
 C(x,y)=\int_\Omega W(x,z)W(y,z)\dd\mu(z),
 \qquad \ell=\TT\tau,
\]
and put $U=1-W$ and
\[
 K_U(x,y)=\int_\Omega U(x,z)U(y,z)\dd\mu(z)\geq0.
\]
The elementary complement-codegree identity is
\begin{equation}
 C(x,y)-r
 =K_U(x,y)+(d(x)-p)+(d(y)-p).
 \label{eq:codegree}
\end{equation}
Indeed, $K_U=1-d(x)-d(y)+C$ and $r=2p-1$.

Fubini and self-adjointness give
\begin{align*}
 H&=\int_{\Omega^2}W(x,y)\ell(x)C(x,y)\dd\mu^2,\\
 F&=\ip\tau A=\ip{\TT\tau}d
   =\int_{\Omega^2}W(x,y)\ell(x)\dd\mu^2.
\end{align*}
If $d=p$ almost everywhere, \eqref{eq:codegree} reduces to
$C-r=K_U\geq0$.  Since $W\geq0$ and $\ell=\TT\tau\geq0$,
\begin{equation}
 H-rF=\int_{\Omega^2}W(x,y)\ell(x)K_U(x,y)\dd\mu^2\geq0.
 \label{eq:regular-repair}
\end{equation}
Finally regularity gives $A=\TT(p\one)=p^2\one$ and hence
$B=\ip dA=p^3$.  Lemma~\ref{lem:anchor} now yields
$H\geq rF\geq r^2B=p^3r^2$.
\end{proof}


%%%%%%%%%%%%%%%%%%%%%%%%%%%%%%%%%%%%%%%%%%%%%%%%%%%%%%%%%%%%%%%%%%%%%%%%%%%%%%%
\section{The exact target condition, a sufficient lemma, and a termwise obstruction}
%%%%%%%%%%%%%%%%%%%%%%%%%%%%%%%%%%%%%%%%%%%%%%%%%%%%%%%%%%%%%%%%%%%%%%%%%%%%%%%

By \eqref{eq:reduction-two}, the original Atlas 130 target is exactly
equivalent to
\begin{equation}
 \ip f{\TT f}+2r(F-rB)+r^2(B-p^3)\geq0.
 \label{eq:exact-target-condition}
\end{equation}
Both correction terms outside the spectral energy are nonnegative by
Lemma~\ref{lem:anchor}.  The following cleaner statement is sufficient for
\eqref{eq:exact-target-condition}, but is potentially strictly stronger.

\begin{conjecture}[Concrete sufficient arbitrary-graphon lemma; unproved]
\label{conj:remaining}
Every graphon of density $p\geq\tfrac12$ satisfies
\begin{equation}
 \boxed{H\geq rF.}
 \label{eq:conjecture}
\end{equation}
\end{conjecture}

This statement is deliberately labelled conjectural and is not asserted to
be equivalent to the original target.  If it holds, then
Lemma~\ref{lem:anchor} immediately gives
$H\geq rF\geq r^2B\geq p^3r^2$, closing Atlas 130 for arbitrary graphons.
The sufficient statement itself is equivalent, by
Proposition~\ref{prop:reduction}, to
\[
 \ip f{\TT f}\geq-r(F-rB).
\]

There is another exact form that exposes a possible complement repair.
For arbitrary $W$, equations \eqref{eq:codegree} and the two Fubini
identities in the proof of Theorem~\ref{thm:regular} give
\begin{align}
 H-rF
 &=\int_\Omega \ell(x)
   \bigl(d(x)^2+A(x)-2p\,d(x)\bigr)\dd\mu(x) \notag\\
 &\quad+\int_{\Omega^2}W(x,y)\ell(x)K_U(x,y)\dd\mu^2.
 \label{eq:complement-repair}
\end{align}
The second term is nonnegative, but the first term is not.

Here is a small exact obstruction to simply discarding the second term.
Take a two-step graphon with block masses
\[
 (m_1,m_2)=\left(\frac{35}{100},\frac{65}{100}\right)
\]
and block matrix
\[
 (W_{ij})=
 \begin{pmatrix}
  2/100&86/100\\
  86/100&41/100
 \end{pmatrix}.
\]
Its edge density is $p=22679/40000>1/2$.  Direct rational evaluation
of the signed first term in \eqref{eq:complement-repair} gives
\begin{equation}
 -\frac{523633393223577}{640000000000000000000}<0,
 \label{eq:signed-obstruction}
\end{equation}
whereas the complete difference is
\begin{equation}
 H-rF=
 \frac{45687431624678040839}{6400000000000000000000}>0.
 \label{eq:repaired-example}
\end{equation}
Thus a proof through \eqref{eq:complement-repair} must retain enough of the
nonnegative complement-codegree term to repair a genuinely negative signed
part.  Neither this example nor the decomposition proves
Conjecture~\ref{conj:remaining}.

There is an even simpler obstruction to a stronger, tempting route.  The
pointwise inequality
\begin{equation}
 \TT\tau\geq rA
 \label{eq:false-pointwise}
\end{equation}
would imply $H\geq rF$ immediately after taking the inner product with
$\tau\geq0$, but \eqref{eq:false-pointwise} is false.  Take the disjoint union
of two clique graphons with block masses $1/10$ and $9/10$, that is,
\[
 (W_{ij})=\begin{pmatrix}1&0\\0&1\end{pmatrix}.
\]
Here $p=41/50$ and $r=16/25$.  On the small block,
\[
 A=\tau=\frac1{100},\qquad
 \TT\tau=\frac1{1000},\qquad
 \TT\tau-rA=-\frac{27}{5000}<0.
\]
This is only a dead-route certificate, not a counterexample to the averaged
conjecture: globally
\[
 F=\frac{1181}{2000},\qquad
 H=\frac{265721}{500000},\qquad
 H-rF=\frac{76761}{500000}>0.
\]
Thus any proof must use averaging or compensation between fibres; it cannot
establish the conjecture by pointwise domination of $\TT\tau$ over $rA$.


%%%%%%%%%%%%%%%%%%%%%%%%%%%%%%%%%%%%%%%%%%%%%%%%%%%%%%%%%%%%%%%%%%%%%%%%%%%%%%%
\section{Exact audit}
%%%%%%%%%%%%%%%%%%%%%%%%%%%%%%%%%%%%%%%%%%%%%%%%%%%%%%%%%%%%%%%%%%%%%%%%%%%%%%%

The companion script
\[
 \texttt{experiments/atlas130\_psd\_regular\_verify.py}
\]
uses only NetworkX, integer arithmetic, and \texttt{fractions.Fraction}.
It verifies the Atlas index, graph6 code and edge set; checks the chromatic
polynomial and target polynomial exactly; recomputes
\eqref{eq:signed-obstruction}--\eqref{eq:repaired-example} from the first
two-step graphon; and checks the exact pointwise obstruction and positive
global gap for the two-clique graphon.  The largest exhaustive loop has only
$8^6$ colour assignments, so the audit has negligible memory requirements.

A second companion,
\[
 \texttt{experiments/atlas130\_conjecture\_search.py},
\]
contains both an exact finite census and a separate numerical falsification
search.  In the exact branch it enumerates all $2^{15}=32768$ symmetric
$5$-by-$5$ binary block matrices, with diagonal entries allowed and all five
blocks of mass $1/5$.  Exactly $16384$ have $p\geq1/2$.  If
$p=P/5^2$, $F=F_0/5^5$, and $H=H_0/5^6$, it evaluates the integer
\[
 5H_0-(2P-25)F_0=5^7(H-rF)
\]
for every matrix.  There are exactly zero negative values, two equality
values, and the least positive gap is exactly
\[
 \frac{128}{5^7}=\frac{128}{78125}.
\]
This is a rigorous certificate for that finite step-graphon family, not a
proof of Conjecture~\ref{conj:remaining}.  The variable-mass continuous search
in the same file remains explicitly numerical evidence only.

\begin{remark}[Status]
Theorems~\ref{thm:psd} and~\ref{thm:regular} are full-interval theorems on
restricted graphon families.  They are not a density interval valid for
all graphons.  Atlas 130 therefore remains open in the arbitrary-graphon
classification, and no green range is claimed here.
\end{remark}

\end{document}
